% DO NOT REMOVE -- This fixes issues where Acrobat on Windows cannot open PDFs
\pdfminorversion=4

% Change to use the correct class file for your paper.
\documentclass[10pt,preprint]{sensys-proc}

%\pagenumbering{arabic}

%\usepackage{adjustbox}    % Auto-resize table content (eg Opo SenSys'14 rel)
\usepackage{amsfonts}     % Adds math fonts, commands such as \begin{align}
\usepackage{array}        % Tables for use in math mode
\usepackage{booktabs}     % Elegant table-formatting library
\usepackage{bold-extra}   % Provides bf+sc (only in textbf+textsc env.)
\usepackage{bytefield}    % Formatting and layout of packets / bytefields
\usepackage[skip=5pt]{caption}
\usepackage[sort,nocompress]{cite} % sort citations
\usepackage{colortbl}     % Color table cells
\usepackage{comment}      % Provides \begin,\end{comment} for large blocks
\usepackage{cprotect}     % Allows verbatim, other formatting in macro args
\usepackage[ampersand]{easylist} % Simpler list syntax
\usepackage{endnotes}     % Footnotes pushed to the end of a document
\usepackage{gensymb}      % Adds useful symbols w/out math mode, e.g. \degree
\usepackage{graphicx}     % For importing graphics
\usepackage{lipsum}       % For generating temporary filler text
\usepackage{listings}     % in-line source code (poorly, consider minted)
\usepackage{mathtools}    % amsmath extension, adds more math formatting
\usepackage[protrusion=true,expansion=true,kerning,spacing]{microtype} % better type, spacing
\usepackage{multirow}     % Multiple row spacing in tables
\usepackage{nth}          % Typeset 33rd correctly as \nth{33}
%\usepackage[section]{placeins} % Don't let figs escape their sections
\usepackage{rotating}     % Rotates any object, note sideways != sidewaysfigure
%\usepackage[all=normal]{savetrees} % For when space is tight, read manual and
                          % selectively enable things. CAN BREAK CONF STYLES!!
\usepackage{soul}         % Provides \hl{} for highlighting
%\usepackage{subfig}
%\usepackage{subfigure}    % Complicated figure creation
\usepackage[subrefformat=parens]{subcaption}   % Replaces both subfig and subfigure
\usepackage{tabularx}     % Complicated table creation
\usepackage{textcomp}     % Provides \textmu for upright mu's
\usepackage{threeparttable} % Add footnotes to a table
%\usepackage{times}        % Force use of the Times font
\usepackage{units}        % For nice fractions, \nicefrac{1}{2} --> 1/2
\usepackage{url}          % Pretty printing of hyperlinks
\usepackage[usenames,dvipsnames,svgnames]{xcolor} % Allow the use and definition of colors
\usepackage{xspace}       % Intelligently add spaces after commands

% The hyperref package must loaded last. Can conflict with some packages, see:
% README ( http://ctan.mackichan.com/macros/latex/contrib/hyperref/README.pdf )
\usepackage[colorlinks=true,citecolor=violet,urlcolor=blue]{hyperref}     % Creates hyperlinks from ref/cite
\hypersetup{pdfstartview=FitH} % Sets default zoom to 100% width

% And, of course, cleveref must be loaded last-last (read: after hyperref)
\usepackage[capitalise,nameinlink,noabbrev]{cleveref}     % Do the right thing with fig/table references

\DeclareCaptionType{copyrightbox}

\newlength\SUBSIZE

\renewcommand{\arraystretch}{1.2} % Space out rows in tables

% Break URLs properly (thanks to Alex Halderman)
\def\UrlBreaks{\do-\do\.\do\@\do\\\do\!\do\_\do\|\do\;\do\>\do\]\do\)\do\,\do\?\do\'\do+\do\=\do\#}
\def\UrlBigBreaks{\do\:\do\/}

% Don't typset URLs in tt font
\urlstyle{same}

%\setlength\paperheight {11in}
%\setlength\paperwidth {8.5in}

% Set the graphics path
\graphicspath{{../figs/}{../images/}}


% No space between bibliography items:
\let\oldthebibliography=\thebibliography
  \let\endoldthebibliography=\endthebibliography
  \renewenvironment{thebibliography}[1]{%
    \begin{oldthebibliography}{#1}%
      \setlength{\parskip}{0ex}%
      \setlength{\itemsep}{0ex}%
  }%
  {%
    \end{oldthebibliography}%
  }

% Tighten up paragraph indents
%\setlength{\parindent}{5mm}

% Some macros that a broadly useful:
\newcommand{\uW}{{\textmu}W\xspace}
\newcommand{\uA}{{\textmu}A\xspace}
\newcommand{\um}{{\textmu}m\xspace}
\newcommand{\us}{{\textmu}s\xspace}
\newcommand{\uF}{{\textmu}F\xspace}
\newcommand{\uJ}{{\textmu}J\xspace}
\newcommand{\iic}{I$^2$C\xspace}

\definecolor{light-gray}{gray}{0.75}
\newcommand{\glipsum}[1][1]{
  \textcolor{light-gray}{\lipsum[#1]}
}

\begin{document}

\title{Hardware and Software for Post-Mote Networked Systems}

\begin{comment}
% AUTHOR STYLE 1
\author{
 \alignauthor{Author 1, Author 2, and Author 3}\\
 \affaddr{Electrical Engineering and Computer Science Department}\\
 \affaddr{University of Michigan}\\
 \affaddr{Ann Arbor, MI 48109}\\
 \email{\{uniqname1,uniqname2,uniqname3\}@umich.edu}
}
\end{comment}

% AUTHOR STYLE 2
\author{
\begin{tabular}{ccc}
  \multicolumn{3}{c}{
    Pat Pannuto$^\dagger$, Bradford Campbell$^\dagger$,
    Michael Anderson$^\ddagger$,
    Amit Levy$^\S$, Tom Bauer$^\S$, Amit Levy$^\S$,
  } \\
  \multicolumn{3}{c}{
    Phillip Levis$^\S$, David Culler$^\ddagger$, and Prabal Dutta$^\dagger$\vspace{0.3cm}
  } \\
  %
  \affaddr{$^\dagger$Computer Science \& Engineering} &
  \affaddr{$^\ddagger$Computer Science \& Engineering} &
  \affaddr{$^\S$Computer Science \& Engineering} \\
  %
  \affaddr{University of Michigan} &
  \affaddr{University of California, Berkeley} &
  \affaddr{Stanford University} \\
  %
  \affaddr{Ann Arbor, MI 48109} &
  \affaddr{Berkeley, CA 94720} &
  \affaddr{Stanford, CA 94305} \\
  %
  \affaddr{\{ppannuto,bradjc,prabal\}@umich.edu} &
  \affaddr{\{m.anderson,culler\}@berkeley.edu} &
  \affaddr{\{levya,tbauer01,pal\}@stanford.edu} \\
\end{tabular}
}

% Michael Andersen <m.andersen@eecs.berkeley.edu> 
% Patrick Pannuto <ppannuto@umich.edu> 
% Brad Campbell <bradjc@umich.edu> 
% Amit Levy <levya@cs.stanford.edu> 
% Tom Bauer <tbauer01@stanford.edu> 
% Philip Levis <pal@cs.stanford.edu> 
% Prabal Dutta <prabal@eecs.umich.edu> 
% David Culler <culler@berkeley.edu> 

\conferenceinfo{Submitted to: SenSys'14}{November 3--6, 2014, Memphis, TN, USA}
\CopyrightYear{2014}
\crdata{XXX-X-XXXXX-XXX-X}

\maketitle

\begin{abstract}
% ABSTRACT

Write abstract here.
Header information is handled by \LaTeX{} magic.


\end{abstract}

% \category{B.4.2}{HARDWARE}{Input/Output and Data
%   Communications}[Input/Output Devices]
% \category{C.3}{COM\-PU\-TER-COM\-MU\-NICATION
%   NET\-WORKS}{Spe\-cial-Purpose and App\-li\-cation-Based Systems}
% \terms{Design, Experimentation, Measurement, Performance}
% \keywords{Mobile phones, Energy harvesting, Phone peripherals, Audio
%   communications, Participatory sensing}

\newpage

%% SUBMISSION INSTRUCTIONS
%% Posters must be submitted as a single PDF containing no more than 3 pages. The
%% first two pages should contain an abstract describing the research content of
%% the poster, along with title, authors, institutional affiliations and contact
%% information. The third page should contain a thumbnail draft of the poster’s
%% contents.

% page limit          % 2
% abstract            % 0.5 pg
\vfill\eject
\section{Introduction}
\label{sec:intro}

Over a decade ago, a flurry of hardware platforms and supporting software
empowered the research community to investigate and explore wireless sensor
networks and their applications. Many papers today still use these
``mote''-class devices, exploring problems in systems software, low-power
networking, and application design.

The underlying technology and its potential applications have progressed a
great deal in the past ten years. The Cortex~M series of 32-bit processors
finally have sleep currents competitive with MSP430 and AVR microcontrollers.
802.15.4 has grown far beyond the destructively closed world of ZigBee, with
new physical layers for new applications. The recent incorporation of
Bluetooth Low Energy (BLE) into mobile phones allows ubiquitous sensing networks to
directly interact with human-centric devices. Also driven by phones, sensors
themselves are orders of magnitude more energy efficient and precise.

Simultaneously, applications have become much richer. Applications in early
sensor network research focused on fixed rate, long-term sensing, guiding a
research agenda of ultra-low power operation and robust multi-hop networking.
Today, personal area networks, tethered to a phone, interact with
proximity networks through iBeacon; intra-vehicle networks connect speakers
and microphones to a phone; building-area networks share knowledge, like
occupancy, among security, HVAC, and lighting monitoring and control.

We have an explosion of new applications, each with new and challenging
requirements. We have reached a turning point in hardware, enabling a whole
new class of device and operational models. It is time for a new platform to
explore and research embedded networked systems in the post-mote era.

         % 0.5 pg
\section{The Post-Mote Landscape}

\begin{easylist}
  & New CPUs mean new OS
  && Low power still critical but tradeoffs and offloads raise interesting problems
  && Not an MCU, not a CPU: how to use novel hardware resources?

  & Perpetual Operation
  && Energy harvesting becoming possible

  & Proximity Networking
  && Not just multihop meshes: BLE, ubiquitious computing
  && Mobility, reversed

  & Privacy Security
  && Network versus system privacy
  && Isolation/privacy domains. Handle address translation, address isolation
\end{easylist}

\glipsum[1-2]
     % 0.5 pg
\section{System Architecture and Principles}

\glipsum[1]

\begin{figure}
  \centering
  \includegraphics[width=\columnwidth]{placeholder}
  \caption{System Architecture}
  \label{fig:system-architecture}
\end{figure}
          % 0.25 pg
\section{Current Status of the Next-Generation Mote}

\begin{easylist}
  & Current status/Demo
  && Storm: Cortex M4
  && BLE stack
  && Perpetual power, hardware/software co-design
\end{easylist}

\glipsum[1]
        % 0.15 pg
%references           % 0.10 pg

%Uncomment this line if your paper has / uses end notes
%\theendnotes

{%\footnotesize
\raggedright
\bibliographystyle{abbrv}
\bibliography{bib}
}

\clearpage
\begin{figure*}
  \centering
  \includegraphics[width=\linewidth]{placeholder}
  \caption{Poster Thumbnail}
  \label{fig:poster-thumbnail}
\end{figure*}


\end{document}

